\documentstyle[]{article}
\begin{document}
\pagestyle{empty}
\samepage
\font\ninept=cmr9
\hsize=10in
\vsize=6.5in
%\hoffset-1.5in
\hoffset-0.75in
\voffset-1.25in
\def\solar{ {\odot} }
\def\lsolar{ {\rm L_{\odot}} }
\def\msolar{ {\rm M_{\odot}} }
\def\mhi{ {\rm M_{\ninept HI}} }
\def\lfir{ {\rm L_{\ninept FIR}} }
\def\lb{ {\rm L_{\ninept B}} }
\def\etal{{\it et~al.}}				% et al. in italics
\def\eg{{{\it e.g.},\ }}			% e.g. in italics
\def\ie{{{\it i.e.},\ }}			% i.e. in italics
\def\cf{{{\it cf.},\ } }			% i.e. in italics
\def\KMS{{$\rm km\ s^{-1}$}}
%
% Slight improvements and fixes to Table 3: July 21, 1993
%
\begin{tabular}{|c|l|c|r|rrr|rrr|}
\multicolumn{10}{c}{{\bf Table~3}.\ 
Derived~Quantitites~for~the~Green~Bank~Sample~of~Polar-Ring~Galaxies}\\ [4pt]\hline
&&&&&&&&&\\
\noalign{\vskip-6pt}
%&PRC 	& Galaxy	& RTele		&Dist_HI&MassHI& L_B	& LFIR & MHI/LB &LFIR/LB &M_HI/LFIR\\
PRC
&\multicolumn{1}{c|}{Galaxy}
&\multicolumn{1}{c|}{Radio}
&\multicolumn{1}{c|}{Distance}
&\multicolumn{1}{c}{$\mhi$}
&\multicolumn{1}{c}{$\lb$}
&\multicolumn{1}{c|}{$\lfir$}
&\multicolumn{1}{c}{$\mhi/\lb$}
&\multicolumn{1}{c}{$\lfir/\lb$}
&\multicolumn{1}{c|}{$\mhi/\lfir$}
\\
No.
&\multicolumn{1}{c|}{Name}
&\multicolumn{1}{c|}{Telescope}
&\multicolumn{1}{c|}{[Mpc]}
&\multicolumn{1}{c}{[$10^9~\msolar$]}
&\multicolumn{1}{c}{[$10^9~\lsolar{\rm _{\ninept ,B}}$]}
&\multicolumn{1}{c|}{[$10^9~\lsolar$]}
&\multicolumn{1}{c}{[$\msolar/\lsolar{\rm _{\ninept ,B}}$]}
&\multicolumn{1}{c}{}
&\multicolumn{1}{c|}{[$\msolar/\lsolar$]}
\\
(1)
&\multicolumn{1}{c|}{(2)}
&\multicolumn{1}{c|}{(3)}
&\multicolumn{1}{c|}{(4)}
&\multicolumn{1}{c}{(5)}
&\multicolumn{1}{c}{(6)}
&\multicolumn{1}{c|}{(7)}
&\multicolumn{1}{c}{(8)}
&\multicolumn{1}{c}{(9)}
&\multicolumn{1}{c|}{(10)}
\\ 
%Table 3.  Derived Quantities
%
% Only FIR values from IRAS FSC2 are used, unless data from IR MAPS or integrated
% scans are available.  Values are those reported by NED.
%
% D-04 (= ESO 296-G11) is a superposition of two galaxies.  HI was detected at two 
% different redshifts, but only total blue and total FIR magnitudes are given.  
% Therefore, the HI masses are computed from the respective redshifts, and the L_B and L_FIR 
% that are shown in the table are computed from each of the two redshifts separately, ie, under 
% the assumption that all of the blue light or all of the FIR flux comes from first one or the 
% other of the two galaxies only.  For this reason, these values and their ratios are only 
% limits.  
%
[4pt]\hline
&&&&&&&&&\\
\noalign{\vskip-6pt}
%
%&PRC 	& Galaxy	& RTele		&Distance&MassHI& L_B	& LFIR  &MHI/LB&LFIR/LB &M_HI/LFIR\\

A-01	& A 0136-0801$^{\ninept 1}$   & GB140		& 75.1 	& 5.07	& 4.52	&       & 1.12	&	&	\\
 	&		& VLA D		& 75.1	& 1.48	& 4.53	&       & 0.32\rlap{6}	&	&	\\
[8pt]
A-02	& ESO 415-G26$^{\ninept 1}$   & GB140		& 60.5	& 5.37	& 7.72	&       & 0.69\rlap{5}	&	&	\\
 	&		& VLA CD/BC	& 60.9	& 3.20	& 7.83	&       & 0.40\rlap{9}	&	&	\\
[8pt]
A-03	& NGC 2685$^{\ninept 2}$      & GB140 	& 13.3	& 1.42	& 4.52	& 0.174& 0.31\rlap{5}	& 0.29\rlap{6}& 8.16$\phantom{0}$  \\
 	&		& WBORK		& 13.2	& 1.21	& 4.47	& 0.173 & 0.27\rlap{1}	& 	& 7.00$\phantom{0}$  \\
 	&		& VLA		& 13.4	& 1.63	& 4.56	& 0.176 & 0.35\rlap{7}	& 	& 9.23$\phantom{0}$  \\
[8pt]
A-05	& NGC 4650A$^{\ninept 1}$     & GB140		& 35.0	& 6.68	& 6.13	& \llap{$<$}1.18$\phantom{0}$ & 1.09	& \llap{$<$}1.48& \llap{$>$}5.66$\phantom{0}$  \\
 	&		& VLA CD/BC	& 35.0	& 4.62	& 6.13	& \llap{$<$}1.18$\phantom{0}$ & 0.75\rlap{3}	& 	& \llap{$>$}3.91$\phantom{0}$  \\
[8pt]
A-06	& UGC 9796$^{\ninept 3}$      & GB140		& 73.8	& 21.0$\phantom{0}$ & 3.37 &  & 6.24 &	&	\\
% 	&		& EFBRG 	& 73.6	& 10.2$\phantom{0}$ & 3.35 &  & 3.04 &	&	\\
 	&		& WBORK		& 73.6	& 4.59	& 3.35	&       & 1.37	&	&	\\
[6pt]\hline 
&&&&&&&&&\\	
\noalign{\vskip-4pt}							        
B-01	& IC 51		& GB140		& 24.2	& 1.59	& 5.75	& 2.40$\phantom{0}$  & 0.27\rlap{7}	& 3.20	& 0.663 \\
B-11	& UGC 5600      & GB140		& 39.6	& 4.45	& 5.19	& 8.59$\phantom{0}$  & 0.85\rlap{7}	& 12.7$\phantom{0}$ & 0.518 \\
B-12	& ESO 503-G17   & GB140		& 134$\phantom{.0}$&$<$23.1$\phantom{0}$& 8.49 &  & $<$2.72&	&	\\
%B-16	& NGC 5122$^{\ninept 3}$      & EFBRG		& 35.4	& 2.07	& 4.99	%&       & 0.41\rlap{4}	&	&	\\
[8pt]
B-17	& UGC 9562      & GB140		& 17.4	& 0.96\rlap{8}& 0.80\rlap{8} &       & 1.20	&	&	\\
B-21	& ESO 603-G21   & GB140		& 43.7	& 6.42	& 1.94	& 4.99$\phantom{0}$  & 3.31	& 19.7$\phantom{0}$	& 1.29$\phantom{0}$  \\
B-23	& A 2330-3751   & GB140		& 	&	& 	&       & 	& 	&	\\
B-27	& ESO 293-IG17  & GB140		& 	&	& 	&       &$<$1.09&	&	\\
[6pt]\hline 
&&&&&&&&&\\ 
\noalign{\vskip-4pt}								        
C-09	& NGC 442       & GB140		& 77.1	& 2.14	& 17.8$\phantom{0}$	&       & 0.12\rlap{1}	& 	&	\\
C-11	& NGC 625       & GB140		&  4.5	& 0.18\rlap{2}& 0.68\rlap{5}& 0.181& 0.26\rlap{6}	& 2.03	& 1.01$\phantom{0}$  \\
C-12	& UGC 1198      & GB140		& 18.7	& 0.13\rlap{5}& 0.93\rlap{7}& 1.31$\phantom{0}$  & 0.14\rlap{4}	& 10.7$\phantom{0}$	& 0.103 \\
C-13	& NGC 660       & GB140		& 13.9	& 8.52	& 5.35	& 21.4$\phantom{00}$  & 1.59	& 29.6$\phantom{0}$	& 0.398 \\
[8pt]
C-14	& NGC 979       & GB140		& 62.4	& 24.1$\phantom{0}$ & 18.7$\phantom{0}$&       & 1.29	&	&	\\
C-18	& ESO 358-G20   & GB140		& 22.9	& 0.21\rlap{7}& 1.25	&       & 0.17\rlap{4}	&	&	\\
C-24	& UGC 4261      & GB140		& 85.9	& 8.02	& 16.6$\phantom{0}$& 12.5$\phantom{00}$& 0.48\rlap{4}	& 5.77	& 0.643 \\
C-25	& UGC 4323      & GB140		& 51.3	& 2.21	& 10.7$\phantom{0}$	&       & 0.20\rlap{7}	& 	&    \\
[4pt]\hline
\end{tabular}
\vfil\eject
\begin{tabular}{|c|l|c|r|rrr|rrr|}
\multicolumn{10}{c}{{\bf Table~3}.\ (continued)}
\\ [4pt]\hline
&&&&&&&&&\\
\noalign{\vskip-6pt}
%&PRC 	& Galaxy	& RTele		&Dist_HI&MassHI& L_B	& LFIR & MHI/LB &LFIR/LB &M_HI/LFIR\\
PRC
&\multicolumn{1}{c|}{Galaxy}
&\multicolumn{1}{c|}{Radio}
&\multicolumn{1}{c|}{Distance}
&\multicolumn{1}{c}{$\mhi$}
&\multicolumn{1}{c}{$\lb$}
&\multicolumn{1}{c|}{$\lfir$}
&\multicolumn{1}{c}{$\mhi/\lb$}
&\multicolumn{1}{c}{$\lfir/\lb$}
&\multicolumn{1}{c|}{$\mhi/\lfir$}
\\
No.
&\multicolumn{1}{c|}{Name}
&\multicolumn{1}{c|}{Telescope}
&\multicolumn{1}{c|}{[Mpc]}
&\multicolumn{1}{c}{[$10^9~\msolar$]}
&\multicolumn{1}{c}{[$10^9~\lsolar{\rm _{\ninept ,B}}$]}
&\multicolumn{1}{c|}{[$10^9~\lsolar$]}
&\multicolumn{1}{c}{[$\msolar/\lsolar{\rm _{\ninept ,B}}$]}
&\multicolumn{1}{c}{}
&\multicolumn{1}{c|}{[$\msolar/\lsolar$]}
\\
(1)
&\multicolumn{1}{c|}{(2)}
&\multicolumn{1}{c|}{(3)}
&\multicolumn{1}{c|}{(4)}
&\multicolumn{1}{c}{(5)}
&\multicolumn{1}{c}{(6)}
&\multicolumn{1}{c|}{(7)}
&\multicolumn{1}{c}{(8)}
&\multicolumn{1}{c}{(9)}
&\multicolumn{1}{c|}{(10)}
\\ 
[4pt]\hline
&&&&&&&&&\\
\noalign{\vskip-6pt}
%&PRC 	& Galaxy	& RTele		&Dist_HI&MassHI& L_B	& LFIR & MHI/LB &LFIR/LB &M_HI/LFIR\\
C-26	& UGC 4332      & GB140		& 76.8	& 5.50	& 12.5$\phantom{0}$& 9.69$\phantom{0}$  & 0.43\rlap{8}	& 5.93	& 0.568 \\
C-27	& UGC 4385      & GB140		& 24.9	& 1.16	& 1.68	& 0.410& 0.69\rlap{1}	& 1.87	& 2.83$\phantom{0}$  \\
C-29	& NGC 2865      & GB140		& 33.5	& 0.78\rlap{1}& 21.4$\phantom{0}$&       & 0.03\rlap{65}	&	&	\\
C-35	& NGC 3414      & GB140		& 17.4	&$<$0.50\rlap{9}& 7.76	& \llap{$<$}0.168&$<$0.06\rlap{57}& \llap{$<$}0.16\rlap{6}&$<$8.67$\phantom{0}$\\
[8pt]
C-38	& NGC 3934      & GB140		& 49.2	& 3.25	& 6.85	& 9.53$\phantom{0}$  & 0.47\rlap{5}	& 10.7$\phantom{0}$	& 0.341 \\
C-39	& NGC 4174      & GB140		& 50.3	& 7.72	& 7.57	&       & 1.02	&	&	\\
C-41	& IC 3370       & GB140		& 39.1	& 0.51\rlap{1}& 52.5$\phantom{0}$& 1.82$\phantom{0}$  & 0.00\rlap{97}	& 0.26\rlap{7}& 0.280 \\
C-42	& NGC 4672      & GB140		& 39.4	& 2.81	& 9.02	& 7.84$\phantom{0}$  & 0.31\rlap{1}	& 6.67	& 0.358 \\
[8pt]
C-44	& NGC 5103      & GB140		& 17.9 & 0.51\rlap{2}& 1.81	&       & 0.28\rlap{2}	&	&	\\
C-45	& NGC 5128  	& GB140		&  3.3	& 0.78\rlap{1}& 20.3$\phantom{0}$& 4.62$\phantom{0}$  & 0.03\rlap{85} & 1.75& 0.169 \\
C-46	& ESO 576-G69   & GB140		& 73.8	& 3.73	& 15.3$\phantom{0}$& 5.39$\phantom{0}$  & 0.24\rlap{4}& 2.71 & 0.693 \\
C-48	& ESO 326-IG6   & GB140		& 111$\phantom{.1}$&$<$20.7$\phantom{0}$& 21.1$\phantom{0}$&       & $<$0.98\rlap{5}&	&	\\
[8pt]
C-50	& UGC 10205     & GB140		& 91.5	& 3.87	& 23.9$\phantom{0}$& 8.43$\phantom{0}$  & 0.16\rlap{2}& 2.70 & 0.459 \\
C-51	& NGC 6285+6286 & GB140		& 77.4	&$<$4.24& 22.6$\phantom{0}$& 112$\phantom{.000}$& $<$0.18\rlap{8}& 38.1$\phantom{0}$ &$<$0.038 \\
C-60	& ESO 464-G31   & GB140		& 87.6	& 9.81	& 24.1$\phantom{0}$	&       & 0.40\rlap{8}	&	&	\\
C-64	& ESO 343-IG13  & GB140		& 76.1	& 6.20	& 13.9$\phantom{0}$& 54.8$\phantom{00}$  & 0.44\rlap{7}	& 30.3$\phantom{0}$	& 0.113 \\
C-69	& NGC 7468      & GB140		& 31.0	& 3.63	& 4.07	& 1.78$\phantom{0}$  & 0.89\rlap{2}	& 3.36	& 2.04$\phantom{0}$  \\
[6pt]\hline 
&&&&&&&&&\\
\noalign{\vskip-4pt}								        
D-02	& NGC 235       & GB140		& 90.2	& 5.45	& 31.2$\phantom{0}$&       & 0.17\rlap{4}	&	&	\\
D-03	& ESO 474-IG28  & GB140		& 88.9	& 4.88	& 14.3$\phantom{0}$& \llap{$<$}3.95$\phantom{0}$  & 0.34\rlap{2}  & \llap{$<$}2.12 & \llap{$>$}1.24$\phantom{0}$  \\
D-04	& ESO 296-G11$^{\ninept 4}$	& GB140		& 68.2	& 15.3$\phantom{0}$& 11.8$\phantom{0}$& 7.07$\phantom{0}$ & 1.30 & 4.60 & 2.17$\phantom{0}$  \\
	&(superposition)& GB140		& 73.5	& 1.76	& 13.7$\phantom{0}$ & 8.21$\phantom{0}$  & 0.12\rlap{8}& 4.60 & 0.214 \\
D-22	& NGC 4643      & GB140		& 11.9	& 0.10\rlap{7}& 4.53	& 0.176& 0.02\rlap{36}	& 0.29\rlap{9}& 0.607 \\
[8pt]
D-25	& UGC 8387      & GB140		& 93.8	& 6.11	& 23.8$\phantom{0}$ & 209$\phantom{.000}$& 0.25\rlap{7}	& 67.3$\phantom{0}$& 0.029 \\
D-28	& NGC 6240      & GB140		& 96.1	& 8.23	& 57.7$\phantom{0}$& 314$\phantom{.000}$& 0.14\rlap{3}	& 41.7$\phantom{0}$	& 0.026 \\
D-30	& ESO 341-IG4   & GB140		& 80.4	& 13.8$\phantom{0}$& 49.5$\phantom{0}$& 13.3$\phantom{00}$  & 0.28\rlap{8} & 2.07& 1.03$\phantom{0}$  \\
D-35	& NGC 7252      & GB140		& 64.3	& 4.46	& 62.0$\phantom{0}$& 28.1$\phantom{00}$  & 0.07\rlap{20} & 3.48 & 0.159 \\
D-43	& ESO 510-G13   & GB140		& 44.4	&$<$3.31& 17.7$\phantom{0}$& 8.09$\phantom{0}$ & $<$0.18\rlap{8}& 3.51	&$<$0.410 \\
[4pt]\hline 
\noalign{\vskip8pt}
\noalign{\vbox{\hsize=7.5in\noindent\baselineskip=12pt

$^{\ninept 1}$\ Results quoted from VLA observations are those of van~Gorkom, 
	Schechter, and Kristian (1987). \hfil\break
$^{\ninept 2}$\ Westerbork observations of NGC~2685 are from Shane (1980); 
	VLA observations are those of Mahon (1992).   \hfil\break
%$^{\ninept 3}$\ Effelsberg observations of UGC~9796 were taken by  
%	W.\ K.\ Huchtmeier (private communication); Westerbork observations  
%\indent	are those of Schechter, \etal\ (1984). \hfil\break
$^{\ninept 3}$\ Westerbork observations of UGC~9796 are those of Schechter, \etal\ (1984). \hfil\break
$^{\ninept 4}$\ ESO~296-G11 is a superposition of two galaxies; we 
	detect HI at two different redshifts.  The quantities $\lb$ and $\lfir$ shown \hfil\break 
\indent in the two rows for this galaxy have been computed as if all of flux 
	originated from one of the two redshifts.  \hfil\break
}}
\end{tabular}
\clearpage
\end{document}



